%=============================================================================
\section{仿真流程}
%=============================================================================

你已经有了离散化的控制器，也有了被控对象的数学模型。在把代码烧进 Flash、拿一台 200 块钱的电机冒险之前，先在电脑上跑一遍不好吗？这一章教你怎么搭仿真，在上硬件之前就干掉 90\% 的 bug。

我们先讲为什么要仿真，再依次走一遍三种仿真测试精度——模型在环（Model-in-the-Loop, MIL）、软件在环（Software-in-the-Loop, SIL）、硬件在环（Hardware-in-the-Loop, HIL）——最后讲怎么往模型里注入真实世界的"脏东西"。文末附了一个直接能用的 Python 模板和一个平衡车案例。

%-----------------------------------------------------------------------------
\subsection{为什么部署前要先仿真}
%-----------------------------------------------------------------------------

硬件调试在比赛季的每个维度上都是昂贵的。打齿轮要花钱，烧电机绕组花钱\textit{还要}等两周快递。就算啥都没坏，每多花一个小时盯着示波器，就是少一个小时去调自瞄或者优化底盘路径规划。

仿真替代不了硬件测试，但它能大幅减少你需要的硬件迭代次数：

\textbf{仿真能告诉你的：}
\begin{itemize}
    \item \textbf{稳定性。}如果闭环极点跑到了 $z$ 平面的右半部分（或者说跑到了单位圆外），那你的增益就是错的。仿真一下就能查出来；硬件查出来的方式是——把齿轮箱给你报废了。
    \item \textbf{大致性能。}上升时间、超调量、稳态时间，模型靠谱的话仿真和实际差个 20--40\%。
    \item \textbf{增益灵敏度。}把某个参数扫一遍 $\pm 50\%$，看控制器还能不能正常工作。如果仿真告诉你惯量加 30\% 就临界稳定了，那你的真实机器人一旦装了块更重的装甲板就会开始振荡。
\end{itemize}

\textbf{仿真不能告诉你的：}
\begin{itemize}
    \item 真实摩擦力（静摩擦、库仑摩擦、温度相关摩擦）。
    \item 减速箱里的间隙——除非你专门建模了。
    \item 传感器噪声的具体特征（量化、漂移、野点）。
    \item 真实通信延迟（CAN 总线仲裁、DMA 抖动）。
\end{itemize}

仿真与现实之间的差距（sim-to-real gap）不可避免。接受它。目标不是完美的仿真，目标是仿真够好——好到你第一次上硬件的结果是"基本能跑，增益微调一下就行"，而不是"电机反转了，把线束直接扯下来了"。

%-----------------------------------------------------------------------------
\subsection{模型在环（MIL）}
%-----------------------------------------------------------------------------

MIL 是最简单也最常见的仿真方式：在同一种语言（Python、MATLAB，甚至 C）里同时写被控对象模型和控制器，跑一个循环，看结果。

\subsubsection{从零搭建直流电机 + PID 仿真}

我们用第三章的简化直流电机模型。忽略电感（当电气时间常数 $\tau_e = L/R$ 远小于控制周期 $T_s$ 时成立），机械动力学为：
\begin{equation}
J \frac{d\omega}{dt} = \frac{K_t}{R}\bigl(V - K_e \omega\bigr) - b\,\omega
\end{equation}
其中 $J$ 是转子惯量，$K_t$ 是力矩常数，$K_e$ 是反电动势常数，$R$ 是绕组电阻，$b$ 是黏性摩擦系数，$V$ 是施加电压，$\omega$ 是角速度。

用前向欧拉离散化（$T_s$ 为采样周期）：
\begin{equation}
\omega[n+1] = \omega[n] + T_s \cdot
    \frac{1}{J}\left[\frac{K_t}{R}\bigl(V[n] - K_e\,\omega[n]\bigr)
    - b\,\omega[n]\right]
\end{equation}

搭配第四章的离散 PID：
\begin{equation}
u[n] = K_p\,e[n] + K_i\,T_s \sum_{k=0}^{n} e[k]
       + \frac{K_d}{T_s}\bigl(e[n] - e[n-1]\bigr)
\end{equation}

下面是完整的 Python 仿真代码：

\begin{verbatim}
import numpy as np
import matplotlib.pyplot as plt

# Plant parameters (small brushed DC motor)
J, b, Kt, Ke, R = 0.01, 0.1, 0.01, 0.01, 1.0
Ts, N = 0.001, 2000          # 1 kHz, 2 seconds
omega = np.zeros(N)
# PID gains
Kp, Ki, Kd = 5.0, 10.0, 0.05
e_int, e_prev = 0.0, 0.0
omega_ref = 100.0             # rad/s step reference

for n in range(N - 1):
    e = omega_ref - omega[n]
    e_int += e * Ts
    e_dot = (e - e_prev) / Ts
    V = Kp * e + Ki * e_int + Kd * e_dot
    V = np.clip(V, -24.0, 24.0)          # actuator saturation
    domega = ((Kt / R) * (V - Ke * omega[n]) - b * omega[n]) / J
    omega[n + 1] = omega[n] + Ts * domega
    e_prev = e

t = np.arange(N) * Ts
plt.plot(t, omega); plt.axhline(omega_ref, ls='--', c='r')
plt.xlabel('Time (s)'); plt.ylabel('omega (rad/s)')
plt.title('DC Motor PID Step Response'); plt.grid(); plt.show()
\end{verbatim}

25 行代码，一个能跑的仿真就有了。跑一下，你会看到阶跃响应有轻微超调，大约 0.5 秒内稳定。如果它发散了，说明增益有问题——在笔记本上发现这事儿，总比在 200 块钱的电机上发现好吧。

\subsubsection{把仿真和理论对照}

跑完阶跃响应之后，直接从仿真数据里量出上升时间 $t_r$、超调量 $M_p$、稳态时间 $t_s$，然后和第三章的二阶近似公式对比：
\begin{align}
M_p &= e^{-\pi\zeta / \sqrt{1 - \zeta^2}} \\
t_s &\approx \frac{4}{\zeta \omega_n}
\end{align}

如果数据在 10--15\% 以内吻合，说明你的仿真和线性化模型是一致的。如果不吻合：
\begin{enumerate}
    \item 仿真可能有 bug（查单位、查正负号）。
    \item 闭环系统不能用二阶来近似（主导极点不是一对干净的共轭复数）。
    \item 执行器饱和把响应削平了，线性分析不再有效。
\end{enumerate}

二阶公式是近似，不是真理。拿不准的时候，信仿真——仿真包含了非线性效应（比如饱和），公式没有。

%-----------------------------------------------------------------------------
\subsection{软件在环（SIL）}
%-----------------------------------------------------------------------------

MIL 适合做设计探索，但有个盲区：你的比赛机器人跑的不是 Python，是 STM32 上的 C，定点运算（或者至少是 32 位浮点），特定的整数类型，ISR 触发的时序。

SIL 的思路很直白：\textbf{把你实际的嵌入式 C 代码编译出来，和 Python 的被控对象模型对接跑。}这能抓住一整类 MIL 抓不到的 bug：
\begin{itemize}
    \item 定点溢出（你的积分项在 65535 处回绕了）。
    \item 整数除法截断（\texttt{int16\_t} 的增益丢精度了）。
    \item 变量类型不匹配（用 \texttt{uint8\_t} 存了一个该用 \texttt{int32\_t} 的值）。
    \item ISR 时序假设在实际中不成立。
\end{itemize}

实现起来很直接。把你的 C 控制器编译成共享库，用 Python 的 \texttt{ctypes} 调用：

\begin{verbatim}
import ctypes

# Compile: gcc -shared -o pid.so -fPIC pid.c
lib = ctypes.CDLL('./pid.so')
lib.pid_update.argtypes = [ctypes.c_float, ctypes.c_float]
lib.pid_update.restype  = ctypes.c_float
lib.pid_init()

for n in range(N - 1):
    V = lib.pid_update(ctypes.c_float(omega_ref),
                       ctypes.c_float(omega[n]))
    # ... plant update same as before ...
\end{verbatim}

现在你仿真里跑的 PID 和你 STM32 上的是字节级别一模一样的。如果积分项在 C 里溢出，在仿真里也一样溢出——你在比赛之前就能抓到。

\textbf{什么时候可以跳过 SIL：}如果你的控制器还是纯 Python 或 MATLAB（比如你正在做快速原型验证，还没写嵌入式代码），MIL 就够了。SIL 只有在你有实际的 C/C++ 代码需要测试时才有价值。

%-----------------------------------------------------------------------------
\subsection{硬件在环（HIL）}
%-----------------------------------------------------------------------------

HIL 更进一步：你的\textbf{真实 MCU} 跑\textbf{真实代码}，但被控对象还是在 PC 上仿真的。PC 和 MCU 通过 UART、SPI 或 CAN 通信——PC 接收 MCU 的控制输出，仿真一步被控对象动力学，再把仿真出来的传感器读数发回去。

\textbf{流程：}
\begin{enumerate}
    \item MCU 启动，正常开始控制循环。
    \item 它不读真实编码器，而是读 PC 通过 UART 发过来的值。
    \item 它计算 PID 输出，把 PWM 占空比（或电压指令）发回 PC。
    \item PC 更新被控对象模型，把新的"传感器"值发回去。
    \item 按控制频率反复循环。
\end{enumerate}

\textbf{什么时候 HIL 值得搞：}
\begin{itemize}
    \item 安全关键系统——控制器 bug 可能伤人。
    \item 非常贵的执行器（工业伺服、液压阀）。
    \item 真实被控对象还不存在（在等定制机加工件）。
\end{itemize}

\textbf{什么时候 HIL 是杀鸡用牛刀：}大多数比赛场景。如果你的电机 30 块钱一个、增益高两倍也不会把齿轮箱炸了，那直接从 MIL 跳到真实硬件就行。你花在搭 HIL 通信协议上的时间，不如直接拿来在真实机器人上调参。

\textbf{延迟警告：}1~kHz 的 HIL 意味着你只有 $\leq 1$~ms 完成 UART 通信 + 被控对象仿真 + 传感器编码的全部往返。在 115200 波特率下，发 8 个字节大约需要 0.7~ms。用更高的波特率（921600 或 USB FS），或者在 HIL 测试时降低控制频率。

%-----------------------------------------------------------------------------
\subsection{让仿真变得真实}
%-----------------------------------------------------------------------------

一个完美的仿真对控制器验证来说毫无意义。如果你的 PID 只有在被控对象模型精确、传感器无噪声时才能跑通，那它在真实硬件上必挂。你必须往仿真里注入和现实一样的"脏东西"。

\subsubsection{传感器噪声}

加入高斯白噪声，方差对齐传感器数据手册（或者测出来的 Allan 方差）。一个典型的 MEMS 陀螺仪噪声密度 $0.01~\text{deg/s}/\sqrt{\text{Hz}}$，在 1~kHz 采样下的离散噪声标准差为 $\sigma = 0.01 \times \sqrt{1000} \approx 0.316$~deg/s。

\subsubsection{扰动}

\begin{itemize}
    \item \textbf{阶跃扰动：}在随机时刻突然加一个负载力矩，测试扰动抑制能力。
    \item \textbf{正弦扰动：}注入 $d(t) = A\sin(2\pi f t)$，在已知频率上测试频域抑制效果。
\end{itemize}

\subsubsection{摩擦}

最简单有用的摩擦模型是把黏性摩擦和库仑摩擦加在一起：
\begin{equation}
\tau_{\text{friction}} = b\,\omega + \tau_c \cdot \operatorname{sign}(\omega)
\end{equation}
其中 $b$ 是黏性系数（我们的电机模型里已经有了），$\tau_c$ 是库仑摩擦力矩。这个模型能抓住零速附近的死区——低速跟踪困难的罪魁祸首。

\subsubsection{通信延迟}

把控制信号延迟 $N_d$ 个采样步来模拟 CAN 总线延迟、传感器处理时间或计算延迟：$V_{\text{applied}}[n] = V_{\text{computed}}[n - N_d]$。

\subsubsection{执行器饱和}

把输出限幅到 $\pm V_{\max}$（或 $\pm \text{PWM}_{\max}$）。这个必须加——线性分析假设控制量无限大，但你的电池是 24~V，电调最大占空比 100\%。

下面演示怎么往仿真循环里加噪声、摩擦和饱和：

\begin{verbatim}
tau_c = 0.005       # Coulomb friction torque (N*m)
sigma = 0.3         # sensor noise std dev (rad/s)
N_delay = 3         # 3-sample control delay
V_buf = [0.0] * (N_delay + 1)

for n in range(N - 1):
    omega_meas = omega[n] + np.random.randn() * sigma    # noisy sensor
    e = omega_ref - omega_meas
    # ... PID computation ...
    V_buf.append(np.clip(V_cmd, -24, 24))                # saturation
    V = V_buf[-(N_delay + 1)]                             # delayed output
    friction = b * omega[n] + tau_c * np.sign(omega[n])   # friction
    domega = ((Kt / R) * (V - Ke * omega[n]) - friction) / J
    omega[n + 1] = omega[n] + Ts * domega
\end{verbatim}

\textbf{核心观点：}一个在完美仿真里能跑但加了噪声和摩擦就挂的控制器，\textit{不是}好控制器。如果你的增益只在天堂里能用，那它们活不过比赛日。

%-----------------------------------------------------------------------------
\subsection{Python 快速仿真模板}
%-----------------------------------------------------------------------------

下面是一个完整的、可以直接复制粘贴的模板。把被控对象参数改成你的机器人的参数，两分钟之内你就有一个能跑的仿真。

\begin{verbatim}
import numpy as np, matplotlib.pyplot as plt

def simulate(plant_fn, Kp, Ki, Kd, ref, Ts=0.001, dur=2.0,
             noise_std=0.0, V_max=24.0):
    N = int(dur / Ts)
    y, u = np.zeros(N), np.zeros(N)
    e_int, e_prev, state = 0.0, 0.0, 0.0
    for n in range(N - 1):
        y_meas = y[n] + np.random.randn() * noise_std
        e = ref - y_meas
        e_int += e * Ts
        u_cmd = Kp * e + Ki * e_int + Kd * (e - e_prev) / Ts
        u[n] = np.clip(u_cmd, -V_max, V_max)
        y[n+1], state = plant_fn(y[n], u[n], state, Ts)
        e_prev = e
    t = np.arange(N) * Ts
    fig, (a1, a2) = plt.subplots(2, 1, sharex=True)
    a1.plot(t, y); a1.axhline(ref, ls='--', c='r'); a1.set_ylabel('Output')
    a2.plot(t, u); a2.set_ylabel('Control'); a2.set_xlabel('Time (s)')
    plt.tight_layout(); plt.show()
    return t, y, u
\end{verbatim}

针对你的具体系统定义一个被控对象函数。对直流电机：

\begin{verbatim}
def dc_motor(omega, V, _, Ts,
             J=0.01, b=0.1, Kt=0.01, Ke=0.01, R=1.0):
    domega = ((Kt/R)*(V - Ke*omega) - b*omega) / J
    return omega + Ts * domega, None

t, y, u = simulate(dc_motor, Kp=5, Ki=10, Kd=0.05,
                    ref=100, noise_std=0.3)
\end{verbatim}

复制过去，改一下被控对象参数，就能跑了。

%-----------------------------------------------------------------------------
\vspace{1em}
\begin{mdframed}[linewidth=1pt, innertopmargin=10pt, innerbottommargin=10pt]
\textbf{案例分析：平衡车部署前在仿真里调 LQR}

\vspace{0.3em}
两轮平衡车就是一个轮子上的倒立摆。状态向量为 $\mathbf{x} = [\theta,\;\dot\theta,\;x,\;\dot x]^T$——机体的倾角和角速度，加上轮子的位置和速度。绕直立平衡点线性化的连续时间动力学为：
\begin{equation}
\dot{\mathbf{x}} = A\mathbf{x} + Bu, \qquad
A = \begin{bmatrix}
0 & 1 & 0 & 0 \\
\frac{g}{l} & -\frac{b_\theta}{J_\theta} & 0 & 0 \\
0 & 0 & 0 & 1 \\
-\frac{g m}{M} & 0 & 0 & -\frac{b_x}{M}
\end{bmatrix}, \quad
B = \begin{bmatrix} 0 \\ -\frac{1}{J_\theta l} \\ 0 \\ \frac{1}{M} \end{bmatrix}
\end{equation}
其中 $g$ 是重力加速度，$l$ 是轮轴到质心的距离，$m$ 是摆体质量，$M$ 是总质量，$b_\theta$、$b_x$ 是摩擦系数。

\textbf{第一步：离散化。}用 ZOH 方法，$T_s = 5$~ms：$A_d = e^{A T_s}$，$B_d = A^{-1}(A_d - I)B$。在 Python 里用 \texttt{scipy.signal.cont2discrete((A, B, C, D), Ts, method='zoh')}。

\textbf{第二步：设计 LQR。}从第六章选代价矩阵 $Q$ 和 $R$。先保守一点：
\[
Q = \operatorname{diag}(100,\; 1,\; 10,\; 1), \qquad R = 1
\]
这意味着倾角的惩罚是角速度的 100 倍，位置的惩罚是速度的 10 倍。求解离散代数 Riccati 方程得到 $K = [K_\theta,\; K_{\dot\theta},\; K_x,\; K_{\dot x}]$。

\textbf{第三步：仿真。}跑闭环系统：$\mathbf{x}[n+1] = A_d \mathbf{x}[n] + B_d(-K\mathbf{x}[n])$，初始条件 $\theta_0 = 5^\circ$。

接下来才是好玩的：\textbf{扫 $Q$ 和 $R$}，探索设计空间。
\begin{itemize}
    \item 加大 $Q_{11}$（倾角惩罚）：倾角恢复更快，控制更激进，轮子加速更大。
    \item 加大 $Q_{33}$（位置惩罚）：机器人漂移更少，但可能牺牲倾角性能。
    \item 加大 $R$：控制更柔和，响应更慢，电流更小。
\end{itemize}

在仿真里扫了 20 分钟之后，队伍定下了 $Q = \operatorname{diag}(500,\; 5,\; 20,\; 1)$，$R = 0.5$，预计稳态时间 0.8~s，超调不超过 $2^\circ$。

\textbf{第四步：部署。}把增益烧进 STM32，开机，推一下。硬件结果：稳态时间 1.1~s，超调约 $2.5^\circ$。30\% 的偏差主要来自减速箱里未建模的摩擦，以及实际质心比假设的稍高一些。但——这才是关键——\textbf{增益第一次上就能用}。没烧电机，没打齿轮，没浪费一下午。在硬件上做一轮微调，稳态时间降到了 0.9~s。

\textbf{教训：}仿真永远不会完全精确。但它能把你送到正确的邻域，这比在硬件上瞎试要值钱得多。
\end{mdframed}
